\scanpage{67}\mspage{(12)}
% archive foliation: 66
\deleted{plus celle d'une}
de ce torseur, \deleted{qui} \inserted{Ce torseur} peut aussi s'interpréter ainsi: la structure d'extension (39) \deleted{comme le torseur des sections} définit la structure d'extension
\[
\tag{43}\diagram{\begin{tikzcd}[ampersand replacement=\&,column sep=large]
\&\underline{\operatorname{End}}(M_0)\arrow[d,phantom,"\rotatebox{90}{$\simeq$}"]\&\&\underline{\operatorname{End}}(M_0)\arrow[d,phantom,"\rotatebox{90}{$\simeq$}"]\&\\
0\arrow[r]\&V_{\underline J}(M_0)\arrow[r]\&V_{\underline J}(M)\arrow[r]\&V_{\underline J}(M_0)\arrow[r]\&1
\end{tikzcd}}
\]
où pour $\mathcal A$ schéma \inserted{(ou faisceau)} en groupes comm. sur $S$, on pose
\[
\tag{44}V_{\underline J}(\mathcal A)\simeq\operatorname{Ker}(\mathcal A^{\underline J}\xrightarrow{\text{somme}}\mathcal A)
\]
(faisceau loc. isom. à $\mathcal A^2$), et $V_{\underline J}$ est un foncteur exact en $\mathcal A$, commutant bien sûr au chgt. de base. Pour $\mathcal A$ annulé par $2$, on a un isom
\[
\tag{45}V_{\underline J}(\mathcal A)\simeq\operatorname{Hom}(M_0,\mathcal A)\qquad\text{si }2\operatorname{id}_{\mathcal A}=0
\]
dû au fait que l'on a
\[
\tag{46}M_0\simeq\mathbb F_2^{\underline J}/\underset{\mathrm{diag}}{\underline{\mathbb F_2}}
\]
\unsure{On a|Donc à} une section \deleted{\unsure{la}} canonique partout non nulle de $V_J(M_0)$, correspondant à l'application d'inclusion $\underline J\unsure{\longleftrightarrow|=}M_0^*\hookrightarrow M_0$, ou encore à la section \unsure{évidente} $\operatorname{id}_{M_0}$ de $\underline{\operatorname{End}}(M_0)$, dont l'image inverse dans (43) est le torseur cherché.
